\documentclass[tikz,border=12pt]{standalone}
\usepackage[utf8]{inputenc}
\usepackage[russian]{babel}
\usepackage{tikz}
\usetikzlibrary{shapes.geometric, arrows.meta, positioning, calc, fit}

\begin{document}

  \begin{tikzpicture}[
    font=\sffamily,
    >=Stealth,
    node distance=1.6cm,
  % Стили блоков
    actor/.style={
      draw=blue!80!black,
      fill=blue!10,
      rounded corners=8pt,
      line width=1.5pt,
      align=center,
      inner sep=10pt,
      text width=8cm
    },
    orchestrator/.style={
      draw=orange!90!black,
      fill=orange!10,
      rounded corners=8pt,
      line width=2pt,
      align=center,
      inner sep=12pt,
      text width=14cm
    },
    storage/.style={
      draw=teal!80!black,
      fill=teal!8,
      rounded corners=6pt,
      line width=1.2pt,
      align=left,
      inner sep=10pt,
      text width=4.3cm,
      minimum height=3.8cm
    },
    consumer/.style={
      draw=purple!80!black,
      fill=purple!10,
      rounded corners=8pt,
      line width=1.5pt,
      align=center,
      inner sep=10pt,
      text width=6.5cm
    },
    arrow/.style={
      ->,
      line width=1.5pt,
      draw=gray!80!black
    },
    lbl/.style={
      font=\small\bfseries,
      fill=white,
      inner sep=2pt
    }
  ]

    % 1. Пользователь (Технолог / Руководитель)
    \node[actor] (user) {
      \textbf{\Large 👤 Технолог / Руководитель фабрики}\\[3pt]
      {\small Заходит в единую страницу Filament: \texttt{/admin/stone/calculator-settings}}
    };

    % 2. Оркестратор (Единый пульт)
    \node[orchestrator, below=of user] (orch) {
      \textbf{\LARGE 🎛️ Единый Оркестратор настроек калькулятора}\\[4pt]
      \textit{\normalsize «Одно окно вместо десятка разрозненных таблиц»}\\[6pt]
      \footnotesize
      \begin{tabular}{cccc}
        \textbf{Вкладка 1:} Формы и лимиты &
        \textbf{Вкладка 2:} Припуски сторон &
        \textbf{Вкладка 3:} Материалы &
        \textbf{Вкладка 4:} Права UI
      \end{tabular}
    };

    % 3. Три системных хранилища под капотом
    \node[storage, below left=2.2cm and 0.2cm of orch] (storeA) {
      \textbf{\large 📐 Формы изделий}\\[-2pt]
      \rule{4.3cm}{0.5pt}\\[3pt]
      \scriptsize
      \textbf{Где:} \texttt{stone\_shapes}\\[3pt]
      • Прямая, Г-, П-образная\\
      • Картинка/чертеж\\
      • Лимиты (длина, ширина)\\
      • Разрешенные вырезы/услуги
    };

    \node[storage, below=2.2cm of orch] (storeB) {
      \textbf{\large 🪚 Профили припусков}\\[-2pt]
      \rule{4.3cm}{0.5pt}\\[3pt]
      \scriptsize
      \textbf{Где:} \texttt{stone\_allowances}\\[3pt]
      • Профиль «Акрил» (+40мм)\\
      • Профиль «Кварц» (+0мм)\\
      • Правила сторон (фасад, торцы, пристенный борт)
    };

    \node[storage, below right=2.2cm and 0.2cm of orch] (storeC) {
      \textbf{\large 💎 Физика слэбов}\\[-2pt]
      \rule{4.3cm}{0.5pt}\\[3pt]
      \scriptsize
      \textbf{Где:} \texttt{ProductType.meta}\\[3pt]
      • Акрил, Кварц, Мрамор...\\
      • Шаг деления (0.25 / 0.5 / 1.0)\\
      • Срез кромки листа (10мм)\\
      • Предел длины (2800мм)
    };

    % 4. Потребители (Витрина и Производство)
    \node[consumer, below left=2.2cm and 0.4cm of storeB] (clientApp) {
      \textbf{\large 💻 Клиентский калькулятор}\\[3pt]
      {\small Показывает клиенту только доступные опции, чертежи и скрывает технические подгибы}
    };

    \node[consumer, below right=2.2cm and 0.4cm of storeB] (cutterApp) {
      \textbf{\large ⚙️ Раскрой и Производство (BOM)}\\[3pt]
      {\small Автоматически увеличивает заготовку на кромку и пилит длинные детали по лимиту доставки}
    };

    % Связи (Стрелки)
    \draw[arrow] (user) -- node[lbl, right=3pt] {Редактирует и жмет «Сохранить»} (orch);

    \draw[arrow] (orch.south) -- ++(0,-0.6) -| node[lbl, near start] {Сохраняет формы} (storeA.north);
    \draw[arrow] (orch.south) -- node[lbl] {Сохраняет припуски} (storeB.north);
    \draw[arrow] (orch.south) -- ++(0,-0.6) -| node[lbl, near start] {Обновляет тип камня} (storeC.north);

    % Связи к потребителям
    \draw[arrow] (storeA.south) -- ++(0,-0.6) -| (clientApp.north);
    \draw[arrow] (storeB.south) -- ++(0,-0.4) -| (clientApp.north);
    \draw[arrow] (storeB.south) -- ++(0,-0.4) -| (cutterApp.north);
    \draw[arrow] (storeC.south) -- ++(0,-0.6) -| (cutterApp.north);

    % Фоновые рамки для пояснения
    \begin{scope}[on background layer]
      \node[draw=gray!40, dashed, fill=gray!3, rounded corners=10pt, fit=(storeA) (storeB) (storeC), inner sep=12pt, label=above:{\small\textbf{База данных (Умные справочники и Типы товаров)}}] {};
    \end{scope}

  \end{tikzpicture}

\end{document}
